\section*{Prologue: A School for Everyone Is Often a School for Nobody}

A founder has decided to build a serious educational institution. Their aspiration is large: a school that eventually teaches, investigates, and contributes to science and technology on an international scale. The founder knows AI, software engineering, statistics, and how to teach. But the institution has no famous alumni, no degree-awarding authority, no campus, no established recruiting channel, no recurring institutional contracts, and little capital.

So the founder announces: \emph{Bayes Institute will teach artificial intelligence to the next generation of engineers.} It sounds consistent with the dream. It also defines a battlefield on which almost every advantage belongs to someone else. Universities have recognition. Large edtech companies have advertising budgets and sales teams. Online platforms have inexpensive content. Creators have audiences. Employers have internal programs. Open-source communities have brilliant free resources.

The founder can respond by doing everything: AI foundations for beginners, MLOps for professionals, data science for analysts, consulting for enterprises, an applied elective for engineering colleges, and perhaps a research seminar. Each offering makes sense individually. Together they demand 5 different promises, 5 kinds of trust, 5 sales motions, and 5 delivery systems. The founder's 20 available hours become 4 hours per business. The institute remains unknown to each customer group.

Now imagine a different choice. The founder identifies 1 small group of experienced developers who have been assigned responsibility for unreliable AI features and need independent, technically demanding feedback. Those developers already possess much of the prerequisite knowledge. Some have a near-term work deadline. They can evaluate a credible sample review. A focused paid studio could serve a few of them without inventing a placement guarantee or a university credential. The founder can speak to real prospective buyers and test whether the pain, willingness to pay, and delivery capacity line up.

This is not yet a business. It may turn out that those engineers prefer employer-sponsored reviews, that their companies prohibit use of work artifacts, that the group is too hard to reach, or that the course format is wrong. But the question is finally tractable. \textbf{The founder has exchanged an unbounded ambition for a falsifiable first position.}

The 2 choices are often confused. A narrow entry market need not imply a narrow eventual institution. The university is the ambition. The first solvable customer problem is the entrance.

Chapters 5 and 6 did 2 different jobs. Chapter 5 showed how to study markets, competition, industry economics, value chains, and opportunity quality. Chapter 6 entered the mind and circumstances of customers: their jobs, triggers, workarounds, budgets, motivations, and purchasing decisions. This chapter asks a harder question that follows from both: \textbf{Given several promising customer groups, where should we concentrate our first scarce resources?}

A founder who gets this wrong can be exceptionally competent and remain invisible. A founder who gets it right can create enough local evidence, trust, learning, and economic momentum to attempt a much larger future.

\subsection*{A Reader's Map}

\begin{table}[htbp]
\centering
\small
\begin{tabularx}{\columnwidth}{l X r}
\toprule
\textbf{Part} & \textbf{Main Question} & \textbf{Sections} \\
\midrule
I & Why begin narrow, and what is a beachhead? & 1--12 \\
II & How do we discover a coherent, underserved group? & 13--27 \\
III & How do we rank entries without being fooled by scores? & 28--43 \\
IV & Is the first market sufficient, reachable, and profitable? & 44--56 \\
V & When should we repeat, expand, or pivot our segment? & 57--70 \\
VI & How does the method change across different industries? & 71--84 \\
VII & How should Bayes Institute choose its first entry? & 85--108 \\
VIII & What should a founder do over the next few weeks? & 109--122 \\
App. & Canvases, interview guides, sizing models, and memos & A--J \\
\bottomrule
\end{tabularx}
\caption{Structural overview of Chapter 7.}
\label{tab:ch7-readers-map}
\end{table}

\subsection*{Evidence Discipline}

This chapter synthesizes economics, strategy, marketing, and entrepreneurship research with practical, explicitly labeled managerial methods. Numerical examples, Bayes-specific market counts, prices, costs, and conversion rates are \textbf{hypothetical teaching assumptions} unless explicitly described as observed. No scorecard predicts startup success with validated precision. Neither a framework nor a famous company's origin story constitutes proof that a particular beachhead will work for Bayes. References are annotated at the end of the chapter.
